\begin{abstract}
Automated fraud detection systems have to make a call on every case, even when the model is not very sure. We introduce \textbf{Safe-L2D-Fraud}, a framework for insurance fraud detection that combines learning to defer (L2D) with Conditional Value-at-Risk (CVaR)-based threshold selection. A pre-trained classifier first generates predictions and uncertainty estimates. Then a deferral policy sends the most unclear claims to human investigators, while also penalising errors in the worst-$\delta$ fraction of cases the system keeps. The deferral decision brings together two signals, Bayesian posterior entropy and calibrated classifier confidence, and uses domain-informed feature selection based on a directed acyclic graph. In a proof-of-concept evaluation on a 1{,}000-claim insurance dataset, Safe-L2D-Fraud reaches 89.5\% system accuracy (95\% CI: [85.0\%, 93.5\%]) at 64\% coverage, improving over standalone XGBoost by 7.5 percentage points, though confidence intervals between methods overlap. We evaluate the framework using bootstrap confidence intervals and cross-validation, and provide a detailed assessment of limitations including small sample size, distribution shift, and adversarial non-stationarity.
\end{abstract}
